\honest{That Tiny Note of Discord}{Eliezer Yudkowsky}{2008}

When I last left my younger self, Eliezer1997 believed that any superintelligence would do what was right, and had an argument for it that I ``fondly claimed to be `formal'.'' \nb{Yudkowsky's 1999 ``Frequently Asked Questions about the Meaning of Life'' does say that the argument ``counted as a formal solution to the Meaning of Life.''} I was on the way to joining the ``smart people who are stupid because they're skilled at defending beliefs they arrived at for unskilled reasons.'' \nb{The line adapts Michael Shermer's ``Smart people believe weird things because they are skilled at defending beliefs they arrived at for non-smart reasons.'' The post does not name Shermer.} But not everyone who digs a hole stays in it, and this is how I got out.

It began with one small thought. By 2000 I thought it proven that building a superintelligence was right if anything was, so no one could have a justifiable conflict of interest over it. That mattered to me, because I found the idea of fighting over the Singularity unbearably stupid; I could not even want to say that my side got there first and would seize the banana. Then: ``Maybe some people would prefer an AI do particular things, such as not kill them, even if life is meaningless?'' The obvious dodge followed from my premises: if life is meaningless, nothing is right, so respecting people's preferences would not be right either. But I did not see myself as someone who dodges; I saw myself as a dutiful rationalist who follows lines of inquiry. So I kept thinking. If people believed they had such preferences, they would have a motive to oppose my project now, and that present conflict would stop right things being done even in the main case, where life has meaning. I could have dismissed the problem with excuses I now hear from others: that it is too hard, from people who think themselves smart enough to build AI; or that there is no time, from people simply not interested.

But I was a perfectionist. My metaethics had seemed to wrap all problems of justice and morality in an airtight wrapper, and this poked a small hole in it, worth patching. To justify the time I thought of Brian Atkins, who paid for the Singularity Institute. He would probably prefer not to die even if life were meaningless, and I did not want to taint our cooperation. English has no simple word for the sentiment; perhaps ``Thou shalt not boil a young goat in its mother's milk.'' Someone who helps you out of altruism should not come to regret it; you owe them that they are really doing what they think they are doing. How would Atkins find out if I did not tell him? I thought that only as the obvious thought of a villain, and I had a counter ready: people are poor deceivers, so I would likely be found out, perhaps by lie detectors invented within thirty years. I lived by the rule that you should be ready to have your thoughts broadcast to the world at any time without embarrassment. These days I allow a self and some secrets, for reasons of Fun Theory; as John McCarthy put it, if everyone lived for others all the time, life would be like a procession of ants following each other in a circle. But on professional topics I still try to be able to pass a future lie detector test, with anyone else willing to take one.

The excuses of time and difficulty were still open. But I did not yet know that the problem would be hard; I had only just thought of it.

So I began to ask: if superintelligences do not produce their own motives from pure logic, how would you specify a fallback morality and write it into the AI? \nb{In 2001, on the SL4 mailing list, the author named a different first step: the possibility that an objective morality ``would have to be built rather than discovered.'' The two accounts may describe different steps of one change.} I already had some safeguards. I knew better than to think an AI needs only One Great Moral Principle; I knew it is wiser to think technologically than politically, that AI programmers should think in concepts that can be written in code, and that suggestively named LISP tokens do not mean anything. These kept me out of some traps that catch other novices. But what mattered was that for the first time I was thinking technically about putting a morality into an AI, without the escape hatch of a mysterious essence of rightness. My earlier philosophizing had never forced me to confront the details; this standard required actual work. Morality slowly became less mysterious, as I began to think inside the black box.

The reasons ``don't matter at all.'' I investigated in detail, so I got better with practice. ``Actions screen off justifications.'' If your arguments justify not working things out in detail, as mine did in 1996, you will not get good at the problem; if they call for detail, you have a chance to build expertise. \nb{This is the moral of ``My Best and Worst Mistake,'' a week earlier, applied to a second case.} This matters because of the ``AI wannabes'' I meet, who have clever reasons not to think about Friendly AI. I name none of them. If your actions do not look good stripped of their justifications, maybe you should reexamine them. Effort will not always save you, since ability can be lacking; but without trying hard you do not even get to the table.

Also, perfectionism really matters. The end of the world does not always come with trumpets and thunder and the highest priority in your inbox. It may first appear as ``one tiny lonely thought,'' which you could dismiss with one easy touch. Understanding then dawned on me slowly, more slowly than it could have. To be continued.
